\documentclass{article}
\usepackage[T1]{fontenc}
\usepackage{lmodern}
\usepackage{graphicx}
\usepackage{amsmath,amssymb}
\usepackage[a4paper,margin=2cm]{geometry}
\usepackage[absolute,overlay]{textpos}
\setlength{\TPHorizModule}{1mm}
\setlength{\TPVertModule}{1mm}
\usepackage[indonesian]{babel}
\usepackage{hyperref}
\setlength{\emergencystretch}{2em}
% unit: Halaman 59

\begin{document}

% === Halaman 59 (python/native) ===

4.  Enkripsi   

• Pesan dapat dienkripsi terlebih dahulu sebelum disembunyikan ke dalam citra. 

• Teknik enkripsi yang sederhana misalnya dengan meng-XOR-kan bit-bit pesan 

dengan bit-bit kunci. Jumlah bit-bit kunci sama dengan jumlah bit pesan.

• Bit-bit kunci dibangkitkan secara acak.

• Kunci untuk pembangkitan bit-bit kunci menjadi stego-key.

• Jika dipakai teknik acak dalam memilih pixel-pixel, maka ada dua stego-key: satu 

untuk pembangkitan bit-bit kunci, satu lagi untuk pembangkitan posisi pixel yang 

dipilih untuk menyembunyikan pesan. 

59

\end{document}
